\begin{claim}{C.74}{Proved}{An exact obstruction to principal-coordinate interpolation}
There is an explicit based loop
$f_\tau\in\mathcal D_\partial$ whose generated group
$G$ is transitive on $\mathbb D$ and whose full
interior centralizer is a closed free proper group
$Q\cong\mathbb R$. The loop has an explicit based
contraction in $\mathcal D_\partial$ commuting with
this actual $Q$.

There are two specified principal coordinates for
the same $Q$-action, with their exact transition
map and inverse. In the second coordinate, write
the loop as
\[
        \widetilde f_\tau(s,t)
                  =(s+u_\tau(t),v_\tau(t)).
\]
The coordinate centralizer homotopy
\[
 (s,t)\longmapsto
       \bigl(s+r u_\tau(t),(1-r)t+r v_\tau(t)\bigr)
             \qquad(0\leq r\leq1)                       \tag{C.74.1}
\]
is a continuous family of interior homeomorphisms,
but its map at $r=1/2,\tau=1/2$ has no continuous
extension at any point of the original boundary
circle. Thus this literal coordinate interpolation
does not give the required relative contraction.

The same map admits no uniform bound by filled
orbits of any compact family of actual
boundary-fixed maps: for every compact
$P\subset\mathcal D_\partial$ and every positive
integer $N$, it fails the inclusion
\[
       A_*(x)\in\operatorname{Fill}\bigl(P^N(x)\bigr)
                         \quad\hbox{for all }x\in\mathbb D,
                                                               \tag{C.74.2}
\]
where $\operatorname{Fill}$ is the actual planar
filled hull of C.70 after the standard disk is
embedded in the plane. The comparison concerns
the stated interpolation and the original maps;
the displayed loop itself remains contractible.
\end{claim}

\begin{proof}
\emph{The actual disk coordinate and three loop paths.}
Use the exact spiral coordinate of C.72, now
writing its real source variables as $(w,t)$:
\[
 \begin{split}
 p(z)&=\frac{z}{1+|z|}
                    \exp\bigl(i\log(1+|z|)\bigr),\\
 p^{-1}(\xi)&=\frac{\xi}{1-|\xi|}
                    \exp\bigl(i\log(1-|\xi|)\bigr),
                    \qquad z=w+it .
 \end{split}                                                   \tag{C.74.3}
\]
These are inverse homeomorphisms between the
real plane and $\mathbb D$ by their polar
formulas, including at zero.

Recall the precise compact-family boundary
control proved in C.72. If $|z'-z|\leq M$ and
$|z|>2M+1$, its direction, radial and logarithm
estimates give
\[
                  |p(z')-p(z)|\leq
                              \frac{4M}{|z|-M}.         \tag{C.74.4}
\]
Consequently every plane homeomorphism with
bounded displacement, and with bounded inverse
displacement, has a unique conjugate extension
fixing $S^1$ pointwise. A compact parameter
family with a common bound and jointly
continuous interior maps and inverses has a
jointly continuous extension with inverses.
The proof is the uniform estimate (C.74.4)
near the circle and the actual interior
conjugacy elsewhere.

Put $\psi(t)=\max\{0,1-|t|\}$ and define
\[
 \begin{split}
 P_{1,a}(w,t)&=(w+a,t),\\
 P_{2,a}(w,t)&=(w,t+a),\\
 P_{3,a}(w,t)&=(w+a\psi(t),t)
                         \qquad(0\leq a\leq1).
 \end{split}                                                   \tag{C.74.5}
\]
Their inverses have the same respective
formulas with $-a$. Their forward and inverse
displacements are at most one. Denote their
actual extensions through (C.74.3) by
$\overline P_{j,a}\in\mathcal D_\partial$.
For $\tau\in[(j-1)/3,j/3]$, set
\[
 \alpha_j(\tau)=
   \min\{6(\tau-(j-1)/3),\,2-6(\tau-(j-1)/3)\},
 \qquad
 f_\tau=\overline P_{j,\alpha_j(\tau)}.                \tag{C.74.6}
\]
Each amplitude is zero at the two endpoints
and one at the midpoint. The three pieces
therefore form one continuous based loop,
with its inverse continuous as well.

There is already an exact relative contraction:
replace each amplitude in (C.74.6) by
$r\alpha_j(\tau)$. The common displacement
bound one still applies. At $r=0$ this is
the identity loop and at $r=1$ it is the
original loop; all joins and the loop
basepoint stay the identity. All three
families commute with horizontal
translations, so this contraction retains
those actual operator equations.

\emph{Transitivity and the full actual interior centralizer.}
The values of the first path generate every
horizontal translation, using products and
inverses of parameters in $[0,1]$.
The values of the second generate every
vertical translation in the same way.
Thus $G$ contains all translations of the
$(w,t)$-plane and is transitive.

Any homeomorphism commuting with all plane
translations is itself a translation:
if its value at zero is $(a,b)$, commutation
forces its value at $(w,t)$ to be
$(w+a,t+b)$. Commutation with $P_{3,1}$
then says
\[
                         \psi(t+b)=\psi(t)
                                      \quad(t\in\mathbb R).
                                                               \tag{C.74.7}
\]
At $t=0$, this gives $\psi(b)=1$, whose unique
solution is $b=0$. Conversely every horizontal
translation commutes with all three families.
Consequently the full interior centralizer,
including possible orientation-reversing
candidates, is exactly
\[
 C_{\operatorname{Homeo}(\mathbb D)}(G|_{\mathbb D})
       =Q=\{p\,q_a\,p^{-1}:a\in\mathbb R\},
       \qquad q_a(w,t)=(w+a,t).                       \tag{C.74.8}
\]
There are no orientation-reversing elements
in this explicitly computed group.

The map $a\mapsto q_a$ is a topological
isomorphism onto its actual interior image:
evaluation at zero gives the inverse, and
the formulas give continuity of the maps
and inverses on every compact set.
This image is closed, since a homeomorphic
limit evaluated at zero has a finite real
parameter and hence is the corresponding
translation. The action is free. It is
proper because two compact source sets
have bounded first coordinates, which
bound every translation parameter taking
one set into the other.
Conjugation by $p$ and its displayed inverse
preserve these assertions. Thus $Q$ has
dimension exactly one and the topology
required in C.33. Its actual extensions to
the closed disk are themselves
boundary-fixed, by (C.74.4).

\emph{Two actual principal sections and their transition.}
The quotient by $Q$ is the real coordinate
$t$. A first section is $\sigma_0(t)=p(0,t)$,
with the principal coordinate
\[
               \Theta_0(w,t)=p(w,t).
\]
A second section of the same quotient is
$\sigma_1(t)=p(-t^3,t)$, and its principal
coordinate is
\[
 \begin{split}
 \Theta_1(s,t)&=p(s-t^3,t),\\
 H=\Theta_0^{-1}\Theta_1 &: (s,t)\longmapsto(s-t^3,t),\\
 H^{-1}&:(w,t)\longmapsto(w+t^3,t).
 \end{split}                                                   \tag{C.74.9}
\]
These are the actual coordinate maps and
their inverse identities. In both,
the same labelled element $p q_a p^{-1}$
acts by addition of $a$ to the first
coordinate. In particular the second
coordinate satisfies every C.33 principal
bundle requirement; no operator has
been changed by the section choice.

At the midpoint $\tau=1/2$ of the second
loop path, the original map is
$pP_{2,1}p^{-1}$. More generally its
vertical-translation path in the second
coordinate is
\[
 F_a(s,t)=
       \bigl(s+(t+a)^3-t^3,t+a\bigr)
             =H^{-1}P_{2,a}H(s,t).                   \tag{C.74.10}
\]
Its literal interpolation with the identity
in that coordinate is
\[
 M_{r,a}(s,t)=
       \bigl(s+r((t+a)^3-t^3),t+ra\bigr).
                                                               \tag{C.74.11}
\]
It is an interior homeomorphism commuting
with every horizontal translation. The
corresponding calculation in the first
coordinate is exact:
\[
 H M_{r,a}H^{-1}(w,t)
  =\bigl(w+3r(1-r)a^2t+r(1-r^2)a^3,\ t+ra\bigr).
                                                               \tag{C.74.12}
\]
In particular $r=1/2,a=1$ gives
\[
 A(w,t)=\left(w+\frac34t+\frac38,t+\frac12\right),
 \qquad
 A^{-1}(w,t)=\left(w-\frac34t,t-\frac12\right).
                                                               \tag{C.74.13}
\]
The actual candidate on the disk is
$A_*=pAp^{-1}$. Formula (C.74.13) is also
an explicit inverse calculation; it
establishes the interior assertion before
any boundary comparison.

\emph{Two sequences at each original boundary point.}
Fix $\zeta=e^{i\theta}\in S^1$ and take
sufficiently large positive integers $n$.
Define
\[
 \begin{split}
 R_n&=e^{2\pi n+\theta}-1,\qquad X_n=(R_n,0),\\
 S_n&=e^{2\pi n+\theta-\pi/2}-1,\qquad Y_n=(0,S_n).
 \end{split}                                                   \tag{C.74.14}
\]
The polar formula (C.74.3) gives
$p(X_n)\to\zeta$ and $p(Y_n)\to\zeta$.
For the horizontal sequence,
$A(X_n)-X_n=(3/8,1/2)$ is bounded, so
(C.74.4) gives $p(A(X_n))\to\zeta$.

For the vertical sequence there is the
exact factorization
\[
        A(Y_n)=\left(S_n+\frac12\right)
                                   \left(\frac34,1\right).
                                                               \tag{C.74.15}
\]
Its argument is $\pi/2-\arctan(3/4)$
and its norm is $\frac54(S_n+\frac12)$.
Thus the logarithmic radial change tends
to $\log(5/4)$, and
\[
       p(A(Y_n))\longrightarrow e^{i\delta}\zeta,
          \qquad
       \delta=\log(5/4)-\arctan(3/4).                  \tag{C.74.16}
\]
The two limits are distinct without any
numerical approximation. The integral
formulas give
$0<\log(5/4)<1/4$ and
$\arctan(3/4)> (3/4)/(1+(3/4)^2)=12/25$.
Consequently
\[
                   -\frac{\pi}{2}<\delta<
                         -\frac{23}{100}<0.           \tag{C.74.17}
\]
In particular $e^{i\delta}\ne1$.
The two sequences (C.74.14) prove that
$A_*$ has no continuous extension at
this specified $\zeta$. Since $\zeta$
was arbitrary, this holds at every
point of the original circle.

\emph{The exact obstruction to the proposed filled-word bound.}
Let $P\subset\mathcal D_\partial$ be any
compact family, and fix $N\geq1$.
The product family $P^N$ is compact.
Joint continuity of evaluation on this
compact family, and its pointwise-fixed
boundary, imply
\[
 \sup_{g\in P^N}|g(p(Y_n))-\zeta|\longrightarrow0.
                                                               \tag{C.74.18}
\]
Every planar filled hull is contained
in the Euclidean convex hull of its
compact set, as proved in C.70.
Therefore every point of
$\operatorname{Fill}(P^N(p(Y_n)))$
also tends uniformly to $\zeta$.
Its proposed member $A_*(p(Y_n))$
instead tends to $e^{i\delta}\zeta$.
This disproves (C.74.2), for these
actual maps and for every such
compact family and finite word length.
Extending each original boundary-fixed
map by the identity outside the closed
disk makes the same assertion in the
original plane homeomorphism group.

The original contraction (C.74.6)
with scaled amplitudes remains valid
and keeps the same operators.
The failed construction is exactly
the interpolation (C.74.1) in the
specified changed principal section.
A section-adapted filling or a
controlled selection still requires
its own boundary argument; this
example does not furnish a nontrivial
centralizer fundamental group.

The executable
\path{code/part_b/c74_principal_interpolation_certificate.py}
checks the coordinate and inverse
identities, the group equations,
the exact interpolation formula,
the two ray calculations and the
rational constants in (C.74.17).
The limit and compact-family
arguments above are topological
proofs, with no floating-point
verification used.
\end{proof}
